% appA.tex -- writer H-walk.  Appendix A: walking the front (proof of the End Theorem).
% Uses the labels of the main text: thm:end, def:chase, lem:F, lem:vflow, lem:seen, lem:chords, lem:excess,
% lem:splittilt, lem:Eprime, lem:duhamel (streamlined labels).  Defines: app:walk, lem:chase, lem:toolkit,
% lem:locality, lem:tiltchain, lem:cont, prop:far, prop:near, prop:length, eq:Emu, eq:Imp.
\providecommand{\eps}{\varepsilon}
\providecommand{\Fam}{\mathcal{F}}
\providecommand{\FamS}{\mathcal{F}_*}
\providecommand{\rhoarc}{\underline{\rho}}
\providecommand{\Tdef}{T_{\mathrm{def}}}

\section{Walking the front: proof of the End Theorem}\label{app:walk}

This appendix proves Theorem~\ref{thm:end}, along the plan sketched after its statement: a continuity argument
(Lemma~\ref{lem:cont}) reduces it to the three estimates of Propositions~\ref{prop:far}, \ref{prop:near} and
\ref{prop:length}; Sections~\ref{app:walk:chase}--\ref{app:walk:tools} collect the bookkeeping these estimates use.

\subsection{Conventions}\label{app:walk:conv}

\textbf{Parameters.} We use the locked numbers of Section~\ref{sec:intro} (in particular $\delta_1=\frac54$),
$U\le U_0=8$, and $\gamma\in\{\frac23,\frac34\}$ (the main
chase uses $\gamma=\frac23$, Stage~A uses $\gamma=\gamma_D=\frac34$; everything below holds for both). As in
Setting~(H), $M=mR^{-\gamma}$ (so $M\lambda^\gamma=m(\lambda/R)^\gamma\le m$ for $\lambda\le R$),
$\eps\le R/K_\eps$, and we assume $R\le1$. The derived constants are
\[
c_6:=\frac{1+\delta_1}{c_A}=\frac92,\qquad
c_\rho^*:=c_A-\frac{2\delta_1}{K_f}=\frac14,\qquad
c_R^*:=c_AK_R-\frac{\delta_1}{K_\eps}=\frac{47}8,\qquad
\chi_g:=\kappa d_L=\frac52,
\]
\[
C_T:=\frac{\sqrt5\,(1-2m^2)^{-1/2}}{1-2^{-\gamma}}\ \ \bigl(<6.11\text{ for }\gamma=\tfrac23,\ <5.58\text{ for }\gamma=\tfrac34\bigr),
\qquad C_\mu:=1+\kappa C_T^2<48 .
\]
Note $c_\rho^*K_f=\frac52\ge2$; this is the inequality needed in Step~1 of the proof of Proposition~\ref{prop:far}.

\textbf{What ``$C$'' means.} A letter $C$ (possibly with an index) is a positive number computed from the locked
parameters only ($m=\frac1{10}$ included). It never depends on $\eps$, $v$, $R$, $\lambda$, $t$, $U$ or on the front;
where $R$ enters, it is written out as a power of $R$. The vision radius enters only in the three ways listed at the
end of Section~\ref{sec:chase} (the third is Lemma~\ref{lem:toolkit}(b) below).

\textbf{How small $v_0$ is.} Always $v\le v_0\le\frac1{10}$. Then, with $\mu=v/(1-v)\le\frac19$,
\begin{equation}\label{eq:V0}
\tfrac1{1-v}\le\tfrac{10}9\le d_L,\qquad e^{2\mu}\le e^{2/9}<1.25=\delta_1 .
\end{equation}
Propositions~\ref{prop:far}--\ref{prop:length} add finitely many conditions $v_0\le R/C$; so $v_0=v_0(m,R)$, as in
Theorem~\ref{thm:end}. Of the number $\delta_t=O_R(v)$ of Setting~(H) we only use $\delta_t\le\frac1{16}$ (true once
$v_0$ is small).

\subsection{The chase}\label{app:walk:chase}

During a leg the pursuer runs straight at a point of the front, and that point flees straight away from it; in the
short time $\eps/(1-v)$ that a leg lasts, the front can stretch only by the factor $e^{2\mu}\approx1$.

\begin{lemma}[Chase basics]\label{lem:chase}
Assume the proviso holds at every non-final leg start that is reached. Then:
\begin{enumerate}
\setlength{\itemsep}{1pt}\setlength{\parskip}{0pt}
\item[(0)] either the chase stops, or $t_j\to\infty$; and $P$ is $1$-Lipschitz;
\item[(1)] during leg $j+1$ the points $x_{a_j},P,x_{a_{j+1}}$ lie on a line in this order, $|x_{a_{j+1}}-P|$
decreases at rate at least $1-v$, and the leg lasts at most $c(a_j,a_{j+1};t_j)/(1-v)\le d_L\eps$;
\item[(2)] every true label $u\in[a_j,a_{j+1}]$ is seen by time $t_j$ (this includes the final leg);
\item[(3)] during leg $j+1$: $l_t(a,b)\le\delta_1l_{t_j}(a,b)$ for all $a<b$, and $|P-x_{a_j}|\le\delta_1\eps$;
\item[(4)] if a path $\tilde P$ equals $P$ on $[0,t_j]$, the working fronts of $P$ and $\tilde P$ agree on $[t_\td,t_j]$.
\end{enumerate}
\end{lemma}

\begin{proof}
By Lemma~\ref{lem:vflow}(a) and $A\le a^*I\le(2/\eps)I$ (Lemma~\ref{lem:F}(ii)), $w_t\le(2v/\eps)w$; hence
$l'\le(2v/\eps)l$ for every pair.

(1) Let $\rho$ solve $\rho'=-1+vf(\rho)$ with $\rho(t_j)=c(a_j,a_{j+1};t_j)>0$. Since $F(\rho\hat u)=f(\rho)\hat u$,
the point $y:=P+\rho\hat u_{j+1}$ solves $y'=vF(y-P)$ with $y(t_j)=x_{a_{j+1}}(t_j)$, so $y=x_{a_{j+1}}$ by uniqueness
while $\rho>0$; and $\rho'\le-(1-v)$ because $f\le1$. Likewise $P-\sigma\hat u_{j+1}$, with $\sigma'=1+vf(\sigma)$,
$\sigma(t_j)=0$, solves the equation of $x_{a_j}$ (as $F(-\sigma\hat u)=-f(\sigma)\hat u$), and $\sigma\ge0$. Finally
$c(a_j,a_{j+1};t_j)\le l_{t_j}(a_j,a_{j+1})\le\eps$ and $1/(1-v)\le d_L$ by \eqref{eq:V0}.

(3) Lengths grow at most by the factor $\exp\bigl((2v/\eps)\cdot\eps/(1-v)\bigr)=e^{2\mu}\le\delta_1$ during the leg.
By (1), $|P-x_{a_j}|=\sigma\le c(a_j,a_{j+1};t)\le l_t(a_j,a_{j+1})\le\delta_1\eps$.
(0) On $[t_\td,T]$, $w\le e^{2v(T-t_\td)/\eps}$; a non-final leg has $l_{t_j}(a_j,a_{j+1})=\eps$ (maximality and
continuity of $l$ in the label), so it moves the base by at least $\eps e^{-2v(T-t_\td)/\eps}$, and only finitely many
legs start before $T$. $P$ is piecewise linear with speed $\le1$.
(2) At $t_j$, $P=x_{a_j}$, so $\rho_u(t_j)=c(a_j,u;t_j)\le l_{t_j}(a_j,u)\le\eps$ by the choice of $a_{j+1}$; apply
Lemma~\ref{lem:seen}. (4) The front at time $t$ depends only on $x^0$ and on $P$ on $[t_\td,t]$ (uniqueness).
\end{proof}

\subsection{The toolkit}\label{app:walk:tools}

\textbf{Reminders from Setting (H).} The base at time $t$ is $a_{j(t)}$, where $j(t)=j$ for $t\in(t_j,t_{j+1}]$ and
$j(t_0)=0$. A leg start $t_j$ ($j\ge1$) belongs to both legs $j$ and $j+1$: a statement about leg $j+1$ may be applied
at $t_j$ with the base $a_j$, since $P(t_j)=x_{a_j}(t_j)$. For a tail pair, $t_W$ is treated as a point of leg
$j_W+1$ (base $a_{j_W}$). \emph{Every pair with $\lambda\le R$ is good}: a pair that is not good has $a<\We$ and
$U-b<\We$, hence $\lambda>U-2\We\ge R+\eps$. Families shrink in time (a pair of the family of leg $k+1$, $k>j$, has
$a\ge a_k\ge a_{j+1}$), and the gap of a fixed label is non-increasing in $t$ and constant on each leg. $\Tdef\le\infty$
is the supremum of the times up to which the chase is defined.

\begin{lemma}[Toolkit]\label{lem:toolkit}
\textup{(CA)} For every good pair with $\lambda\ge\eps/2$ and every $t\ge t_\td$: $c(a,b;t)\ge c_A\min(\lambda,K_RR)$.
In particular $c(a,b;t)\ge c_A\lambda$ for every pair with $\eps/2\le\lambda\le R$.

\textup{(c) (Distance from the base.)} Let $t$ lie in leg $j+1$ and let $s\ge a_j$ be a label with $U-s\ge\We$ or
$a_j\ge\We$. Then
\[
\rho_s(t)\ge c_A\min(s-a_j,K_RR)-\delta_1\eps,\qquad \max(\rho_s(t),\eps)\ge\min(s-a_j,K_RR)/c_6 .
\]
Now let $t'\le\Tdef$ and assume that \textup{(J1)} holds at every leg start $t_i\le t'$ for all pairs of $\FamS(t_i)$.
\begin{enumerate}
\setlength{\itemsep}{1pt}\setlength{\parskip}{0pt}
\item[\textup{(a)}] \emph{(Legs.)} Let $t_j\le t'$ be a leg start. (a1) If $U-a_j\ge\eps/2$ and leg $j+1$ is not final,
then $a_{j+1}-a_j\in[\eps/\kappa,\eps]$ and $l_{t_j}(a_j,a_{j+1})=\eps$. (a2) If $U-a_j<\eps/2$, leg $j+1$ is final; if
$a_j<\We$, it is not final. (a3) If leg $j+1$ is not final, $c(a_j,a_{j+1};t_j)\ge c_A(a_{j+1}-a_j)>0$: the proviso
holds. (a4) At most $\kappa U/\eps$ non-final legs start in $[t_\td,t']$, each lasts at most $d_L\eps$, and
$t'-t_\td\le\chi_gU\le20$.
\item[\textup{(b)}] \emph{(Gap--time integrals.)} Fix a label $a$ and a set $\mathcal T\subset[t_\td,t']$ of times at
which $g(\tau):=a-a_{j(\tau)}\ge0$. The values of $g$ on distinct legs differ by at least $\eps/\kappa$, and:
(i) if $g_{\min}$ is the smallest value of $g$ on $\mathcal T$ and $h\ge0$ is non-increasing on $[g_{\min},\infty)$,
then $\int_{\mathcal T}h(g(\tau))\,d\tau\le d_L\eps\,h(g_{\min})+\chi_g\int_{g_{\min}}^\infty h(y)\,dy$;
(ii) if $g\le\bar g$ on $\mathcal T$ and $h\ge0$ is non-decreasing on $[0,\bar g+\eps/\kappa]$, then
$\int_{\mathcal T}h(g(\tau))\,d\tau\le\chi_g\int_0^{\bar g+\eps/\kappa}h(y)\,dy$;
(iii) at most $\kappa x/\eps+1$ legs meet $\{\tau\in\mathcal T:g(\tau)\le x\}$.
\end{enumerate}
\end{lemma}

The point of (b): a leg lasts about $\eps$ and the gap drops by about $\eps$ per leg, so ``time at gap $g$'' is
comparable to ``$dg$''.

\begin{proof}
(CA) Chords are non-decreasing in time (Lemma~\ref{lem:chords}(a)); apply (S1). Pairs with $\eps/2\le\lambda\le R$ are
good, and $R\le K_RR$.

(c) Put $\Delta:=s-a_j$ and $y:=\min(\Delta,K_RR)$. If $\Delta<\eps/2$, the first right-hand side is negative and the second is
$<\eps$. Otherwise $(a_j,s)$ is good, and by (CA) and Lemma~\ref{lem:chase}(3),
$\rho_s\ge|x_s-x_{a_j}|-|P-x_{a_j}|\ge c_Ay-\delta_1\eps$. If $c_Ay-\delta_1\eps\ge y/c_6$ we are done; otherwise,
since $c_A-1/c_6=c_A\delta_1/(1+\delta_1)$, we get $y<(1+\delta_1)\eps/c_A=c_6\eps$, so $\max(\rho_s,\eps)\ge\eps>y/c_6$.

(a) Always $a_{j+1}-a_j\le l_{t_j}(a_j,a_{j+1})\le\eps$, since $w\ge1$. The pairs $(a_j,s)$, $\eps/2\le s-a_j\le\eps$,
lie in $\Fam(t_j)$ and are admissible (if $t_j<t_W$, then $a_j<\We$ and $U-s\ge U-\We-\eps\ge\We$), so (J1) applies.
(a1) The label $s:=a_j+\eps/2\le U$ has $l_{t_j}(a_j,s)\le\kappa\eps/2<\eps$, so $a_{j+1}\ge s$. If the leg is not
final, $l_{t_j}(a_j,a_{j+1})=\eps$ by maximality, and (J1) for $(a_j,a_{j+1})$ gives $\eps\le\kappa(a_{j+1}-a_j)$.
(a2) A final leg has $U-a_j\le l_{t_j}(a_j,U)\le\eps$, so $a_j\ge U-\eps>\We$ if $\We>0$; thus legs with $a_j<\We$ are
not final. Now let $U-a_j<\eps/2$ (so $a_j\ge\We$ and $t_j\ge t_W$); we argue by induction on $j$. For $j=0$,
$l_{t_0}(0,U)=U<\eps/2$. For $j\ge1$, leg $j$ was walked, so by induction $U-a_{j-1}\ge\eps/2$ and by (a1)
$\lambda_1:=a_j-a_{j-1}\in[\eps/\kappa,\eps]$ with $l_{t_{j-1}}(a_{j-1},a_j)=\eps$. The pair $(a_{j-1},U)$ has length
$\lambda_1+(U-a_j)\in[\eps/2,\frac32\eps]\subset[\eps/2,R]$ and lies in $\Fam(t_j)=\FamS(t_j)$, so by (J1) and the
monotonicity of lengths in $t$,
$l_{t_j}(a_j,U)\le\kappa\cdot\frac32\eps-\eps=\frac78\eps<\eps$, i.e.\ $a_{j+1}=U$. (Here $\kappa<\frac43$ is used.)
(a3) By (a1),(a2), $\eps/2\le a_{j+1}-a_j\le\eps\le R$; apply (CA).
(a4) Each non-final leg moves the base by at least $\eps/\kappa$ (by (a1),(a2)) and the moves add up to at most $U$; each
lasts at most $d_L\eps$ by Lemma~\ref{lem:chase}(1). Hence $t'-t_\td\le(\kappa U/\eps)d_L\eps=\chi_gU\le\frac52\cdot8$.

(b) Non-final legs move the base by at least $\eps/\kappa$ (part (a)), so the values of $g$ are spaced by
$\eps/\kappa$; the final leg occupies no time. (i) Let $g_1>\dots>g_n=g_{\min}$ be the values. For $i<n$,
$g_i-\eps/\kappa\ge g_{i+1}\ge g_{\min}$ and $h$ is non-increasing, so $h(g_i)\le(\kappa/\eps)\int_{g_i-\eps/\kappa}^{g_i}h$;
these intervals are disjoint. Each leg lasts at most $d_L\eps$, so the integral is at most
$d_L\eps\,h(g_{\min})+d_L\eps(\kappa/\eps)\int_{g_{\min}}^\infty h$. (ii) Likewise
$h(g_i)\le(\kappa/\eps)\int_{g_i}^{g_i+\eps/\kappa}h$ for every value, with disjoint intervals in
$[0,\bar g+\eps/\kappa]$. (iii) The values in $[0,x]$ are spaced by $\eps/\kappa$.
\end{proof}

\textbf{Start times and bases.} For a pair $(a,b)\in\FamS(t)$ its \emph{start time} is $t_\circ:=t_\td$ if it is a body
pair and $t_\circ:=t_W$ if it is a tail pair; at its start time a pair satisfies (S2) resp.\ (TW). For
$\tau\in[t_\circ,t]$ let $\bar a(\tau):=a_{j(\tau)}$ and $\bar a^+(\tau):=a_{j(\tau)+1}$ be the base and the next base,
with the convention $\bar a(t_W):=a_{j_W}$, $\bar a^+(t_W):=a_{j_W+1}$ for tail pairs.

The next lemma says that the bootstrap bounds may be used, at all times since the start time, on the pair itself and
on the pairs lying between the pursuer and it; the only subtlety is the excused tail.

\begin{lemma}[Locality]\label{lem:locality}
Let $t\le\Tdef$, $(a,b)\in\FamS(t)$ and $\tau\in[t_\circ,t]$; write $\bar a=\bar a(\tau)$, $\bar a^+=\bar a^+(\tau)$.
\begin{enumerate}
\setlength{\itemsep}{1pt}\setlength{\parskip}{0pt}
\item[(i)] $(a,b)$ is a pair of the family of the leg of $\tau$ and belongs to $\FamS(\tau)$; its gap
$g(\tau):=a-\bar a\ge0$ is non-increasing on $[t_\circ,t]$ and constant on each leg; and $a\ge\bar a^+$ if $g(\tau)>0$.
\item[(ii)] Every pair of the family of the leg of $\tau$ contained in $[\bar a,\max(b,\bar a^+)]$ belongs to
$\FamS(\tau)$.
\item[(iii)] Every $s\in[a,b]$ satisfies the hypothesis of Lemma~\ref{lem:toolkit}(c) at $\tau$ with the base $\bar a$.
\end{enumerate}
\end{lemma}

\begin{proof}
(i) $a\ge a_{j(t)}\ge\bar a$. If $\tau$ lies in the leg of $t$, the family is the same; otherwise
$a\ge a_{j(t)}\ge a_{j(\tau)+1}=\bar a^+$, so $(a,b)$ is in the family of the leg of $\tau$. For a tail pair at
$\tau=t_W$, (TW) gives $a-a_{j_W}\ge K_f\lambda\ge5\eps>a_{j_W+1}-a_{j_W}$, so again $a\ge\bar a^+$. Body pairs are
admissible at all times, tail pairs for $\tau\ge t_W$. The bases are non-decreasing, and $g(\tau)>0$ means $a\ne\bar a$,
hence $a\ge\bar a^+$.
(ii) If $\tau\ge t_W$, then $\FamS(\tau)=\Fam(\tau)$, which contains the family of the leg of $\tau$ (at $t_W$ the
family of leg $j_W$ contains that of leg $j_W+1$). If $\tau<t_W$, then $(a,b)$ is a body pair and $\bar a<\We$, so
$\bar a^+\le\bar a+\eps<\We+\eps$; a pair $(a'',b'')$ in $[\bar a,\max(b,\bar a^+)]$ has
$U-b''\ge\min(U-b,\,U-\We-\eps)\ge\We$ (recall $U\ge2\We+R+\eps$ if $\We>0$), so it is a body pair.
(iii) For a body pair $U-s\ge U-b\ge\We$; for a tail pair $\tau\ge t_W$ and $\bar a\ge a_{j_W}\ge\We$.
\end{proof}

Two nested chords whose lengths differ by at most a factor $4$ are nearly parallel, because the excess of the longer
one limits how far the shorter one can tilt (Lemma~\ref{lem:splittilt}). A general nested pair is joined by a chain of
intervals that halve in length, and the tilts add up to a geometric series.

\begin{lemma}[Tilt chain]\label{lem:tiltchain}
Let $t$ lie in leg $j+1$, let $I=(a,b)$ be a pair of the family of leg $j+1$, and suppose that \textup{(J1)},
\textup{(J2)} hold at time $t$ for every pair of that family contained in $[a_j,\max(b,a_{j+1})]$. Then for every pair
$J\subset I$ of that family,
\[
|e_J-e_I|\le C_TM\lambda_I^\gamma=C_Tm(\lambda_I/R)^\gamma ,
\]
and for $a<s\le b$ and $\lambda_s:=\max(s-a,\eps/2)$, with $e:=e_{ab}$,
\begin{equation}\label{eq:Emu}
\int_a^s|T-e|^2\,ds\le M^2\lambda^{2\gamma}\min(C_\mu\lambda_s,\lambda),\qquad
\int_a^s|T-e|\,ds\le M\lambda^\gamma\bigl(\kappa\lambda_s\min(C_\mu\lambda_s,\lambda)\bigr)^{1/2}.
\end{equation}
\end{lemma}

\begin{proof}
All pairs below lie in $[a_j,b]$ and in the family, so (J1), (J2) and (CA) apply to them.

\emph{Step 1 (one step).} Let $J\subsetneq I$ with $\rho:=\lambda_I/\lambda_J\le r\le4$. By
Lemma~\ref{lem:splittilt}(i), $Z^2\le N(I)^2\,\ell(1-\ell/L)$ with $\ell:=l(J)\le\kappa\lambda_J$ and
$L:=l(I)\le\kappa\lambda_I$. The function $\ell(1-\ell/L)$ increases with $L$, so
$Z^2\le N(I)^2\max_{0<\ell\le\kappa\lambda_J}\ell(1-\ell/(\kappa\lambda_I))=N(I)^2\kappa\lambda_J\mathsf q(\rho)$, where
$\mathsf q(\rho)=\rho/4$ for $\rho\le2$ (interior maximum) and $\mathsf q(\rho)=1-1/\rho$ for $\rho\ge2$ (endpoint).
With (J2) and $c(J)\ge c_A\lambda_J$,
\[
\frac Z{c(J)}\le\frac{M\lambda_I^{1/2+\gamma}(\kappa\lambda_J\mathsf q(\rho))^{1/2}}{c_A\lambda_J}
=\frac{\kappa^{1/2}}{c_A}\,\mathsf s(\rho)\,M\lambda_I^\gamma,\qquad \mathsf s(\rho):=(\rho\mathsf q(\rho))^{1/2}
=\begin{cases}\rho/2,&\rho\le2,\\(\rho-1)^{1/2},&\rho\ge2,\end{cases}
\]
and $\kappa^{1/2}/c_A=\sqrt5$; $\mathsf s$ is increasing, $\mathsf s(2)=1$, $\mathsf s(4)=\sqrt3$. Moreover, as
$M\lambda_I^\gamma\le m$, we have $N(I)^2/(4c(J))\le M^2\lambda_I^{1+2\gamma}/(4c_A\lambda_J)\le m^2r/2\le2m^2<1$. So
Lemma~\ref{lem:splittilt}(ii) gives $|e_J-e_I|\le\sqrt5\,c_m\,\mathsf s(r)M\lambda_I^\gamma$ with
$c_m:=(1-2m^2)^{-1/2}$.

\emph{Step 2 (the chain).} Let $J=(a',b')$, $\lambda_k:=2^{-k}\lambda_I$ and $K:=\min\{k\ge1:\lambda_k\le\lambda_J\}$.
\emph{Case A: $a\ge a_{j+1}$, or $a'=a$, or $\lambda_I/2\ge\eps$.} For $0\le k<K$ put $\alpha_k:=\min(a',b-\lambda_k)$,
$I_k:=(\alpha_k,\alpha_k+\lambda_k)$, and $I_K:=J$. Then $I_0=I$; $J\subset I_k\subset I$ (as $\alpha_k\in[a,a']$ and the
right end $\min(a'+\lambda_k,b)$ lies in $[b',b]$); $I_{k+1}\subset I_k$ ($\alpha_k$ is non-decreasing and the right
end non-increasing in $k$). Each $I_k$ is in the family: its length lies in $(\lambda_J,\lambda_I]\subset[\eps/2,R]$,
and its first point is either $a'$ (the first point of the family pair $J$) or $b-\lambda_k\ge a+\lambda_I/2$ with
$k\ge1$; the latter is $\ge a_{j+1}$, trivially if $a\ge a_{j+1}$, and because $a+\eps\ge a_j+\eps\ge a_{j+1}$ if
$\lambda_I/2\ge\eps$; if $a'=a$ then $\alpha_k=a$ throughout. Consecutive length ratios are $2$, except the last one,
which lies in $[1,2]$. Step~1 with $r=2$ gives $|e_{I_{k+1}}-e_{I_k}|\le\sqrt5\,c_m M\lambda_k^\gamma$, and summing
the geometric series,
$|e_J-e_I|\le\sqrt5\,c_m M\lambda_I^\gamma\sum_{k\ge0}2^{-k\gamma}=C_TM\lambda_I^\gamma$.
\emph{Case B: otherwise}, i.e.\ $a=a_j<a'$ and $\lambda_I<2\eps$. Then $\lambda_I/\lambda_J<2\eps/(\eps/2)=4$, and
Step~1 with $r=4$ gives $|e_J-e_I|\le\sqrt{15}\,c_m M\lambda_I^\gamma\le C_TM\lambda_I^\gamma$ (as
$\sqrt3<(1-2^{-\gamma})^{-1}$).

\emph{Step 3 (\eqref{eq:Emu}).} $J:=(a,a+\lambda_s)\subset I$ is a family pair (first point $a$, length in
$[\eps/2,\lambda]$). By Lemma~\ref{lem:excess}(b) with $\hat u=e$, (J2), (J1) and Step~2,
$\int_a^s|T-e|^2\le\int_J|T-e|^2=N(J)^2+c(J)|e_J-e|^2\le M^2\lambda_s^{1+2\gamma}+\kappa\lambda_sC_T^2M^2\lambda^{2\gamma}
\le C_\mu M^2\lambda^{2\gamma}\lambda_s$; also $\int_a^s|T-e|^2\le\int_I|T-e|^2=N(I)^2\le M^2\lambda^{1+2\gamma}$. The
second bound follows by Cauchy--Schwarz, since $l(a,s)\le l(J)\le\kappa\lambda_s$.
\end{proof}

\subsection{The continuity scheme}\label{app:walk:cont}

\begin{lemma}[Continuity scheme]\label{lem:cont}
Put $\bar\theta:=\frac58$. Suppose that whenever $t^*\le\Tdef$ is such that the chase has not stopped before $t^*$ and
\textup{(J1)}, \textup{(J2)} hold for all $t\in[t_\td,t^*]$ and all pairs of $\FamS(t)$, then for all
$t\in[t_\td,t^*]$ and all pairs of $\FamS(t)$
\begin{equation}\label{eq:Imp}
\textup{(Imp)}\qquad l_t(a,b)\le\tfrac98\lambda\qquad\text{and}\qquad E(a,b;t)\le\tfrac12\bar\theta M^2\lambda^{1+2\gamma}.
\end{equation}
Then the proviso holds at every non-final leg start, the chase stops, \textup{(J1)}, \textup{(J2)} hold for all
$t\in[t_\td,T_\stopp]$ and all pairs of $\FamS(t)$, and $t_W$ is reached if $\We>0$.
\end{lemma}

\begin{proof}
Let $\tau^*$ be the supremum of the $t'\le\Tdef$ such that the chase has not stopped before $t'$ and (J1), (J2) hold on
$[t_\td,t']$. At $t_\td$ we have $l=\lambda$ and, by (S2), $N\le s_0M\lambda^{1/2+\gamma}$ (all pairs of
$\FamS(t_\td)$ are body pairs), so $\tau^*\ge t_\td$.

\emph{The invariants hold at $\tau^*$} (if $\tau^*<\infty$). If $\tau^*\in(t_j,t_{j+1}]$, the family is constant on
$(t_j,\tau^*]$, and so is $\FamS$ unless $\tau^*=t_W$; as $l,E$ are continuous in $t$ (Lemma~\ref{lem:vflow}(c)), (J1),
(J2) pass to the limit. If $\tau^*=t_W$, the newly admitted pairs are tail pairs, covered by (TW). So (Imp) holds on
$[t_\td,\tau^*]$.

\emph{Extension.} Suppose the chase does not stop at $\tau^*$. If $\tau^*=t_j$ is a leg start, leg $j+1$ is not final
and the proviso holds (Lemma~\ref{lem:toolkit}(a3) with $t'=\tau^*$), so the chase is defined on
$[\tau^*,\tau^*+\delta_0]$ for some $\delta_0>0$ with no leg start and no $t_W$ in $(\tau^*,\tau^*+\delta_0]$; there
$\FamS(t)\subset\FamS(\tau^*)$. For a fixed pair, $l'\le(2v/\eps)l$ and $|E'|=v|\int(T-e)\cdot AT|\le(4v/\eps)l$
(Lemmas~\ref{lem:Eprime}(i), \ref{lem:F}(ii)). Choose $\delta\in(0,\delta_0]$ with $\frac98e^{2v\delta/\eps}<\kappa$ and
$\delta(4v/\eps)\kappa\le\frac12(1-\bar\theta)M^2(\eps/2)^{2\gamma}$. Then for $t\in[\tau^*,\tau^*+\delta]$ and every
pair of $\FamS(t)$ (recall $\lambda\ge\eps/2$): $l_t\le\frac98e^{2v\delta/\eps}\lambda<\kappa\lambda$, and
$E(t)\le\frac12\bar\theta M^2\lambda^{1+2\gamma}+\delta(4v/\eps)\kappa\lambda\le\frac12M^2\lambda^{1+2\gamma}$. So
(J1), (J2) hold on $[t_\td,\tau^*+\delta]$, contradicting the definition of $\tau^*$.

Hence the chase stops at $\tau^*$, unless $\tau^*=\infty$; but then it would have infinitely many legs (by
Lemma~\ref{lem:chase}(0)), against Lemma~\ref{lem:toolkit}(a4). So $\tau^*=T_\stopp<\infty$, the proviso held at every
non-final leg start (Lemma~\ref{lem:toolkit}(a3)), and (J1), (J2) hold on $[t_\td,T_\stopp]$. Finally, if $\We>0$, the final leg $J+1$ has $U-a_J\le l_{t_J}(a_J,U)\le\eps$, so
$a_J\ge U-\eps>\We$ and $j_W\le J$.
\end{proof}

\textbf{Standing assumptions for Sections~\ref{app:walk:far}--\ref{app:walk:length}.} Let $t^*$ be as in
Lemma~\ref{lem:cont}: the chase has not stopped before $t^*$, and (J1), (J2) hold for all $t\in[t_\td,t^*]$ and all
pairs of $\FamS(t)$. Then Lemma~\ref{lem:toolkit}(a),(b) hold with $t'=t^*$. Fix $t\in[t_\td,t^*]$ and
$(a,b)\in\FamS(t)$, put $e:=e_{ab}$, and use the notation $t_\circ,\bar a,\bar a^+,g$ of Lemma~\ref{lem:locality};
$\hat g:=\min(g(t),K_RR)$. For a tail pair at $t=t_W$, (TW) gives (Imp) directly ($l\le\frac{17}{16}\lambda$,
$E\le\frac12s_0^2M^2\lambda^{1+2\gamma}$), so for tail pairs we take $t>t_W$. By Lemma~\ref{lem:locality}, at every
$\tau\in[t_\circ,t]$: the pairs of the family of the leg of $\tau$ in $[\bar a,\max(b,\bar a^+)]$ satisfy (J1), (J2) and
Lemma~\ref{lem:tiltchain}; Lemma~\ref{lem:toolkit}(c) applies at every $s\in[a,b]$ with base $\bar a$; and $P$ lies on
$[x_{\bar a},x_{\bar a^+}]$ (Lemma~\ref{lem:chase}(1)). For $s\in[a,b]$ we write $\Delta_s:=s-\bar a(\tau)\ge g(\tau)$, $\hat\rho_s:=\max(\rho_s,\eps)$, and
$\rhoarc:=\min_{[a,b]}\rho_s$; $p_a:=(x_a-P)\times e$, so $|p_a|$ is the distance from $P$ to the line through
$x_a,x_b$. Note $l\ge\lambda$ (as $w\ge1$) and, by (J1), $l\le\kappa\lambda$.

\subsection{Far pairs}\label{app:walk:far}

\begin{proposition}[Far pairs]\label{prop:far}
Let $v_0\le\min\bigl(\frac1{300},\frac R{10},\frac{mR}{C_1}\bigr)$ with the number $C_1$ of \eqref{eq:farfinal}. If
$(a,b)$ is far at $t$, then it is far at every $\tau\in[t_\circ,t]$, and $N(a,b;t)\le f_FM\lambda^{1/2+\gamma}$.
Moreover, at every $\tau\in[t_\circ,t]$, with $g=g(\tau)$:
\begin{enumerate}
\setlength{\itemsep}{1pt}\setlength{\parskip}{0pt}
\item[(i)] $\rhoarc\ge\eps$ and $\rhoarc\ge2l$; $\rhoarc\ge c_\rho^*g=g/4$ if $g\le K_RR$, and
$\rhoarc\ge c_R^*R=\frac{47}8R$ if $g>K_RR$;
\item[(ii)] if $g\le c_{\mathrm{off}}R$: $|p_a|\le C_pMg^{1+\gamma}$, where $C_p:=(\kappa C_\mu)^{1/2}(1+1/K_f)^{1/2+\gamma}<9$.
\end{enumerate}
\end{proposition}

\begin{proof}
The gap is non-increasing, so the pair is far on $[t_\circ,t]$. Fix $\tau\in[t_\circ,t]$ and $g=g(\tau)$; then
$g\ge K_f\lambda\ge K_f\eps/2=5\eps$.

\emph{Step 1: distance (this is where $\delta_1=\frac54$ is needed).} For $s\in[a,b]$, $\Delta_s\ge g$, so
Lemma~\ref{lem:toolkit}(c) gives $\rho_s\ge c_A\min(g,K_RR)-\delta_1\eps$. If $g\le K_RR$: since
$\eps\le2\lambda\le2g/K_f$, $\rhoarc\ge(c_A-2\delta_1/K_f)g=g/4$, and $g/4\ge5\eps/4\ge\eps$ --- at the smallest far gap
$g=5\eps$ this is exactly the inequality $c_\rho^*K_f/2\ge1$; also $g/4\ge K_f\lambda/4=\frac52\lambda\ge2l$. If
$g>K_RR$: as $\eps\le R/K_\eps$, $\rhoarc\ge c_AK_RR-\delta_1R/K_\eps=\frac{47}8R$, which is $\ge\eps$ and
$\ge\frac{47}8\lambda\ge4l$. This proves (i). Hence Lemma~\ref{lem:duhamel} applies on $[t_\circ,t]$
($c(a,b;t_\circ)>0$ by (CA)), with $\mathcal A\le(1+\frac1{2\pi}+\frac2{4\pi^2})v/\rhoarc\le\frac54v/\rhoarc$.

\emph{Step 2: amplification.} We claim, for $t_\circ\le\tau\le t$,
\begin{equation}\label{eq:amp}
\exp\int_\tau^t\mathcal A\le A_0\Bigl(\frac{\min(g(\tau),K_RR)}{\hat g}\Bigr)^{\nu}.
\end{equation}
The times with $g>K_RR$ form an initial segment of total length at most $20$ (Lemma~\ref{lem:toolkit}(a4)), on which
$\mathcal A\le\frac54\cdot\frac8{47}\frac vR<\frac v{4R}$; they contribute at most $5v/R$. On the times of $[\tau,t]$ with
$g\le K_RR$ (present only if $g(t)\le K_RR$, and then $\hat g=g(t)$) we have $\mathcal A\le5v/g$, and
Lemma~\ref{lem:toolkit}(b)(i) with $h(y)=y^{-1}\one\{y\le\min(g(\tau),K_RR)\}$ and $g_{\min}=g(t)\ge5\eps$ gives
\[
\int\mathcal A\le5v\Bigl[\frac{d_L\eps}{g(t)}+\chi_g\log\frac{\min(g(\tau),K_RR)}{g(t)}\Bigr]
\le2v+\frac{25}2v\log\frac{\min(g(\tau),K_RR)}{\hat g}.
\]
Since $v\le\frac1{300}$ and $v\le R/10$: $5v/R+2v<\log2=\log A_0$ and $\frac{25}2v\le\frac1{24}=\nu$; as the ratio is
$\ge1$, \eqref{eq:amp} follows.

\emph{Step 3: forcing.} Write $\mathcal B:=2vl^{3/2}|p_a|/(\pi\rhoarc^3)$ for the forcing term of
Lemma~\ref{lem:duhamel}.
\emph{Case $g>c_{\mathrm{off}}R$.} At the label $s$ where $\rho_s=\rhoarc$, $|p_a|\le\rho_a\le\rhoarc+|x_a-x_s|\le\rhoarc+l\le\frac32\rhoarc$,
so $\mathcal B\le\frac3\pi\kappa^{3/2}v\lambda^{3/2}/\rhoarc^2\le2v\lambda^{3/2}/\rhoarc^2$. On $\{g>K_RR\}$
(duration $\le20$, $\rhoarc>5R$) this integrates to at most $\frac85v\lambda^{3/2}/R^2$. On $\{c_{\mathrm{off}}R<g\le K_RR\}$,
$\rhoarc\ge g/4$, and Lemma~\ref{lem:toolkit}(b)(i) with $h(y)=y^{-2}$, $g_{\min}>R/2$ and $d_L\eps\le R/5$ gives
$\int g^{-2}d\tau\le\frac R5\cdot\frac4{R^2}+\chi_g\frac2R<\frac6R$, hence a contribution at most
$32\cdot6\,v\lambda^{3/2}/R$. Altogether, using $R\le1$,
\begin{equation}\label{eq:forcefar}
\int_{\{g>c_{\mathrm{off}}R\}}\mathcal B\,d\tau\le\frac{194\,v\lambda^{3/2}}{R^2}.
\end{equation}
\emph{Case $g\le c_{\mathrm{off}}R$: the offset.} $P(\tau)$ lies on $[x_{\bar a},x_{\bar a^+}]$ and the distance to a line is convex,
so $|p_a|\le\max_{a'\in\{\bar a,\bar a^+\}}Y_{a'}$ with $Y_{a'}:=|(x_a-x_{a'})\times e|$. Fix $a'$, put $g':=a-a'\in[g-\eps,g]$,
$L:=(a',a)$ and $I':=(a',b)$. Both are pairs of the family of the leg (first point $a'$; lengths $g'\ge4\eps$ and
$g'+\lambda\le(1+\frac1{K_f})g\le\frac{11}{20}R<R$), contained in $[\bar a,\max(b,\bar a^+)]$, so (J1), (J2) and
Lemma~\ref{lem:tiltchain} apply to them; note $(a,b)\subset I'$. Since $(x_a-x_{a'})\times e=\int_L(T-e)\times e\,ds$,
Cauchy--Schwarz and Lemma~\ref{lem:excess}(b) give
$Y_{a'}\le(l(L)\,V_e(L))^{1/2}$, and
\[
V_e(L)\le V_e(I')=N(I')^2+c(I')|e_{I'}-e|^2
\le M^2\Lambda^{1+2\gamma}+\kappa\Lambda\,C_T^2M^2\Lambda^{2\gamma}=C_\mu M^2\Lambda^{1+2\gamma},
\]
with $\Lambda:=g'+\lambda\le(1+\frac1{K_f})g$ and $l(L)\le\kappa g$. Thus $Y_{a'}\le C_pMg^{1+\gamma}$, which is (ii).
With $\rhoarc\ge g/4$ and $l\le\kappa\lambda$:
\begin{equation}\label{eq:forcenear}
\mathcal B\le C_{\mathcal B}\,vM\lambda^{3/2}g^{\gamma-2}\quad\text{on }\{g\le c_{\mathrm{off}}R\},\qquad C_{\mathcal B}:=\frac{128\kappa^{3/2}C_p}{\pi}<513.
\end{equation}

\emph{Step 4: result.} By Lemma~\ref{lem:duhamel} on $[t_\circ,t]$ and \eqref{eq:amp} (on $\{g\le c_{\mathrm{off}}R\}$,
$\min(g(\tau),K_RR)=g(\tau)$ and $\hat g=g(t)$),
\[
N(t)\le A_0\Bigl(\frac{\min(g(t_\circ),K_RR)}{\hat g}\Bigr)^{\!\nu}N(t_\circ)
+A_0\Bigl(\frac{K_RR}{\hat g}\Bigr)^{\!\nu}\int_{\{g>c_{\mathrm{off}}R\}}\mathcal B\,d\tau
+A_0\int_{\{g\le c_{\mathrm{off}}R\}}\Bigl(\frac{g(\tau)}{g(t)}\Bigr)^{\!\nu}\mathcal B(\tau)\,d\tau .
\]
\emph{First term.} By (S2) for a body pair (where $g(t_\circ)=a$), or by (TW) for a tail pair (where
$g(t_\circ)=g_W$), $N(t_\circ)\le s_0M\lambda^{1/2+\gamma}\min(1,K_f\lambda/G_\circ)^\nu$ with $G_\circ:=\min(g(t_\circ),K_RR)$. So the
weight of $s_0M\lambda^{1/2+\gamma}$ is $(\min(G_\circ,K_f\lambda)/\hat g)^\nu$; this is $\le(K_f\lambda/g(t))^\nu\le1$ if
$g(t)\le K_RR$, and $\le1$ otherwise (then $\hat g=K_RR\ge G_\circ$). So the term is at most
$A_0s_0M\lambda^{1/2+\gamma}=\frac14M\lambda^{1/2+\gamma}$.
\emph{Second term.} $(K_RR/\hat g)^\nu\le(K_R/K_f)^\nu(R/\lambda)^\nu\le2(R/\lambda)^\nu$, and since $\gamma+\nu<1$ and
$\lambda\le R$, $\lambda^{3/2}(R/\lambda)^\nu\le\lambda^{1/2+\gamma}R^{1-\gamma}$. With \eqref{eq:forcefar} and
$R^{-1-\gamma}=M/(mR)$ the term is at most $776\,(v/mR)\,M\lambda^{1/2+\gamma}$.
\emph{Third term.} By \eqref{eq:forcenear} it is at most $2C_{\mathcal B} vM\lambda^{3/2}g(t)^{-\nu}\int h(g(\tau))\,d\tau$
with $h(y):=y^{\nu+\gamma-2}$, non-increasing because $\nu+\gamma<1$. Lemma~\ref{lem:toolkit}(b)(i) ($g_{\min}=g(t)$)
and $\eps\le2\lambda$, $g(t)\ge\lambda$ give
\[
g(t)^{-\nu}\!\int h\le d_L\eps\,g(t)^{\gamma-2}+\frac{\chi_g\,g(t)^{\gamma-1}}{1-\gamma-\nu}
\le\Bigl(2d_L+\frac{\chi_g}{1-\gamma-\nu}\Bigr)\lambda^{\gamma-1}\le16\,\lambda^{\gamma-1}
\]
($1-\gamma-\nu\ge\frac5{24}$). So the term is at most $32C_{\mathcal B} vM\lambda^{1/2+\gamma}$.
Adding up, with $m\le1$, $R\le1$:
\begin{equation}\label{eq:farfinal}
N(t)\le\Bigl(\frac14+\frac{C_1v}{4mR}\Bigr)M\lambda^{1/2+\gamma},\qquad C_1:=4(776+32C_{\mathcal B})<7\cdot10^4,
\end{equation}
which is $\le\frac12M\lambda^{1/2+\gamma}=f_FM\lambda^{1/2+\gamma}$ because $v\le mR/C_1$.
\end{proof}

\subsection{Near pairs}\label{app:walk:near}

During the near phase (of duration of order $\lambda$) the excess grows at most at rate $v\,m_\lambda^2\,H(g/\lambda)$,
where $H$ has an integrable singularity at the pursuer. The singular kernels are integrated with the following
elementary inequality; the constant $C_*$ produced is finite and independent of $\eps$, which is all we need.

\begin{lemma}[Layer cake]\label{lem:layercake}
Let $\varpi$ be a positive measure on $[0,1]$ with $\varpi([0,\sigma])\le\Phi(\sigma)$ for all $\sigma\in[0,1]$, where
$\Phi$ is $C^1$ and non-decreasing, and let $\psi\ge0$ be non-increasing. Then
$\int\psi\,d\varpi\le\psi(0)\Phi(0)+\int_0^1\psi\Phi'\,d\sigma$.
\end{lemma}

\begin{proof}
Write $\int\psi\,d\varpi=\int_0^\infty\varpi(\{\psi>u\})\,du$. Each $\{\psi>u\}$ is an interval $[0,\sigma_u)$ or
$[0,\sigma_u]$, whose $\varpi$-measure is at most $\Phi(\sigma_u)$, which is the measure of the same set under
$\Phi(0)\delta_0+\Phi'd\sigma$.
\end{proof}

\begin{proposition}[Near pairs]\label{prop:near}
There is a number $C_*$ (depending only on the locked parameters; given in \eqref{eq:Cstar}) such that if $(a,b)$ is
near at $t$, then
\[
E(a,b;t)\le\bigl(\tfrac12f_F^2+C_*v\bigr)M^2\lambda^{1+2\gamma}.
\]
In particular, if $v_0\le3/(16C_*)$, then $E(a,b;t)\le\frac12\bar\theta M^2\lambda^{1+2\gamma}$ with
$\bar\theta=\frac12(1+f_F^2)=\frac58$.
\end{proposition}

\begin{proof}
Put $m_\lambda:=M\lambda^\gamma\le m$. We measure everything in units of the pair length $\lambda$: at a time $\tau$ of
a leg with gap $g$ and for $s\in[a,b]$, put
\[
\vartheta:=\frac\eps\lambda\in(0,2],\qquad
\zeta:=\frac g\lambda\in[0,K_f),\qquad
\sigma:=\frac{s-a}\lambda\in[0,1]
\]
(the normalised vision radius, gap, and position in the pair), so that $\Delta_s=\lambda(\zeta+\sigma)$.

\emph{Start value.} The gap is constant on legs and non-increasing, so the times of $[t_\circ,t]$ at which the pair is
near form an interval $(t_n,t]$, where $t_n$ is a leg start at which the pair is far, or $[t_\td,t]$ (then $t_\circ=t_\td$,
since tail pairs are far at $t_W$; put $t_n:=t_\td$). In the second case $E(t_n)\le\frac12s_0^2M^2\lambda^{1+2\gamma}$ by
(S2); in the first, Proposition~\ref{prop:far} at time $t_n$ (or (TW) if $t_n=t_W$) gives
$E(t_n)\le\frac12f_F^2M^2\lambda^{1+2\gamma}$. On $(t_n,t]$, Lemma~\ref{lem:Eprime}(iii) with $\bar q=0$ gives
$E'\le\mathrm I+\mathrm{II}$, where $\mathrm I:=v\int_a^ba^*(\rho_s)|T-e|^2\,ds$ and $\mathrm{II}:=v\int_a^b|\varsigma||q|\,ds$.

\emph{Envelope.} Since $\Delta_s<(K_f+1)\lambda\le11R<K_RR$, Lemma~\ref{lem:toolkit}(c) gives $\hat\rho_s\ge\max(\eps,\Delta_s/c_6)
=\lambda\max(\vartheta,(\zeta+\sigma)/c_6)\ge\lambda(\zeta+\sigma+\vartheta)/(2c_6)$. From Definition~\ref{def:F},
$a^*(\rho_s)\le2/\hat\rho_s$. Hence
\begin{equation}\label{eq:env}
a^*(\rho_s)\le\frac{4c_6}{\lambda(\zeta+\sigma+\vartheta)},\qquad\frac1{\hat\rho_s}\le\frac{2c_6}{\lambda(\zeta+\sigma+\vartheta)} .
\end{equation}
\emph{Budgets.} Let $\Omega(\sigma):=\int_a^{a+\lambda\sigma}|T-e|^2ds$ and $G(\sigma):=\int_a^{a+\lambda\sigma}|\varsigma|\,ds$.
By \eqref{eq:Emu} with $\lambda_s=\lambda\max(\sigma,\vartheta/2)\le\lambda(\sigma+\vartheta)$, and $|\varsigma|\le|T-e|$,
\begin{equation}\label{eq:budget}
\Omega(\sigma)\le C_\mu m_\lambda^2\lambda(\sigma+\vartheta),\qquad G(\sigma)\le(\kappa C_\mu)^{1/2}m_\lambda\lambda(\sigma+\vartheta).
\end{equation}
\emph{Term I.} By \eqref{eq:env}, \eqref{eq:budget} and Lemma~\ref{lem:layercake} with $\Phi=C_\mu m_\lambda^2\lambda(\sigma+\vartheta)$:
\[
\mathrm I\le4c_6C_\mu vm_\lambda^2\Bigl(\frac\vartheta{\zeta+\vartheta}+\int_0^1\frac{d\sigma}{\zeta+\sigma+\vartheta}\Bigr)
\le4c_6C_\mu vm_\lambda^2H_1(\zeta),\qquad H_1(\zeta):=1+\log\Bigl(1+\frac1{\zeta+\vartheta}\Bigr).
\]
\emph{Term II at long-type times} ($g+\lambda>R$). By Lemma~\ref{lem:F}(iii), $|q|\le\rho_s/\hat\rho_s^2\le1/\hat\rho_s$, so by
\eqref{eq:env}, \eqref{eq:budget}, $\mathrm{II}\le2c_6(\kappa C_\mu)^{1/2}vm_\lambda H_1(\zeta)$. Here
$R/\lambda<\zeta+1<K_f+1$, so $m_\lambda=m(\lambda/R)^\gamma>m(K_f+1)^{-\gamma}$, and therefore
$\mathrm{II}\le2c_6(\kappa C_\mu)^{1/2}(K_f+1)^\gamma m^{-1}\,vm_\lambda^2H_1(\zeta)$.

\emph{Term II at short-type times} ($g+\lambda\le R$). By Lemma~\ref{lem:F}(iii), $|q(s)|\le|(x_s-P)\times e|/\hat\rho_s^2$,
and $|(x_s-P)\times e|\le|(x_s-x_a)\times e|+|p_a|\le G(\sigma)+|p_a|$ (as $|T\times e|=|\varsigma|$). This splits
$\mathrm{II}\le\mathrm{II}_a+\mathrm{II}_b$. For $\mathrm{II}_a=v\int G|\varsigma|\hat\rho_s^{-2}ds$ we use the measure
$d(G^2/2)$, with $G^2/2\le\frac12\kappa C_\mu m_\lambda^2\lambda^2(\sigma+\vartheta)^2$, and the kernel $4c_6^2/(\lambda(\zeta+\sigma+\vartheta))^2$:
\[
\mathrm{II}_a\le2c_6^2\kappa C_\mu vm_\lambda^2\Bigl(\frac{\vartheta^2}{(\zeta+\vartheta)^2}+\int_0^1\frac{2(\sigma+\vartheta)\,d\sigma}{(\zeta+\sigma+\vartheta)^2}\Bigr)
\le4c_6^2\kappa C_\mu vm_\lambda^2H_1(\zeta).
\]
For $\mathrm{II}_b=v|p_a|\int|\varsigma|\hat\rho_s^{-2}ds$ we use $dG$:
\[
\mathrm{II}_b\le\frac{4c_6^2(\kappa C_\mu)^{1/2}vm_\lambda|p_a|}{\lambda}\Bigl(\frac\vartheta{(\zeta+\vartheta)^2}+\int_0^1\frac{d\sigma}{(\zeta+\sigma+\vartheta)^2}\Bigr)
\le\frac{8c_6^2(\kappa C_\mu)^{1/2}vm_\lambda|p_a|}{\lambda(\zeta+\vartheta)} .
\]
\emph{The offset.} As $P\in[x_{\bar a},x_{\bar a^+}]$, $|p_a|\le\max(Y_{\bar a},Y_{\bar a^+})$, $Y_{a'}:=|(x_a-x_{a'})\times e|$.
\begin{itemize}
\setlength{\itemsep}{1pt}\setlength{\parskip}{0pt}
\item $g>0$. Then $a\ge\bar a^+$ (Lemma~\ref{lem:locality}(i)) and $g\ge\bar a^+-\bar a\ge\eps/\kappa>\eps/2$. For
$a'\in\{\bar a,\bar a^+\}$ with $g':=a-a'>0$: if $g'\ge\eps/2$, $(a',a)$ is a family pair and $l(a',a)\le\kappa g'$; if
$g'<\eps/2$, then $a'=\bar a^+$ and $(a',a)\subset(\bar a^+,\bar a^++\eps/2)$, a family pair, so $l(a',a)\le\kappa\eps/2$.
In both cases $l(a',a)\le\kappa\lambda(\zeta+\vartheta)$. The pair $I':=(a',b)$ is a family pair (length
$\Lambda\le g+\lambda\le R$) in $[\bar a,\max(b,\bar a^+)]$, and as in Step~3 of Proposition~\ref{prop:far},
$\int_{I'}|T-e|^2\le C_\mu M^2\Lambda^{1+2\gamma}$ with $\Lambda\le(1+K_f)\lambda$. By Cauchy--Schwarz,
$Y_{a'}\le(l(a',a)\int_{I'}|T-e|^2)^{1/2}\le(\kappa C_\mu)^{1/2}(1+K_f)^{1/2+\gamma}m_\lambda\lambda(\zeta+\vartheta)^{1/2}$;
so $\mathrm{II}_b\le8c_6^2\kappa C_\mu(1+K_f)^{1/2+\gamma}vm_\lambda^2H_2(\zeta)$ with $H_2(\zeta):=(\zeta+\vartheta)^{-1/2}$.
\item $g=0$. Then $a=\bar a$ and $Y_{\bar a}=0$; if the leg is final, $P=x_{\bar a}$ and $p_a=0$. Otherwise: if
$\bar a^+\le b$, $Y_{\bar a^+}\le\int_a^b|T-e|\le(l\,N^2)^{1/2}\le\kappa^{1/2}m_\lambda\lambda$; if $\bar a^+>b$ (so
$\lambda<\eps$, $\vartheta>1$), $I_0:=(\bar a,\bar a^+)$ is a family pair of length $\le\eps$ containing $(a,b)$, and by
Lemma~\ref{lem:tiltchain}, $Y_{\bar a^+}=c(I_0)|e_{I_0}\times e|\le\kappa\eps\,C_TM\eps^\gamma\le2^{1+\gamma}\kappa C_Tm_\lambda\lambda$.
So $\mathrm{II}_b\le32c_6^2(\kappa C_\mu)^{1/2}\kappa C_T\,vm_\lambda^2/\vartheta$ (using $2^{1+\gamma}\le4$ and
$\kappa^{1/2}\le2\kappa C_T$).
\end{itemize}
\emph{Per-leg rate.} Let $C_2$ be the largest of the numbers
$4c_6C_\mu+2c_6(\kappa C_\mu)^{1/2}(K_f+1)^\gamma/m$, $\ 4c_6C_\mu+4c_6^2\kappa C_\mu+8c_6^2\kappa C_\mu(1+K_f)^{1/2+\gamma}$,
$\ 32c_6^2(\kappa C_\mu)^{1/2}\kappa C_T$. We have shown: on a leg with gap $\zeta>0$,
$E'\le C_2vm_\lambda^2\bigl(H_1(\zeta)+H_2(\zeta)\bigr)$; on the leg with $g=0$, $E'\le C_2vm_\lambda^2(H_1(0)+1/\vartheta)$.

\emph{Sum over the near legs.} Each leg lasts at most $d_L\eps=d_L\vartheta\lambda$. At most one leg has $g=0$; it
contributes at most $d_L\vartheta\lambda\,C_2vm_\lambda^2(H_1(0)+1/\vartheta)\le4d_LC_2vm_\lambda^2\lambda$, since
$\vartheta\le2$ and $\vartheta\log(1+1/\vartheta)\le1$. The positive gaps $\zeta_i$ are distinct, $\ge\vartheta/\kappa$, spaced by at
least $\vartheta/\kappa$, and $<K_f$; for the non-increasing $H:=H_1+H_2$ this gives
$\sum_iH(\zeta_i)\le H(\vartheta/\kappa)+(\kappa/\vartheta)\int_0^{K_f}H$, so these legs contribute at most
$d_LC_2vm_\lambda^2\lambda\bigl(\vartheta H(\vartheta/\kappa)+\kappa\int_0^{K_f}H\bigr)$. Here
$\vartheta H_1(\vartheta/\kappa)\le\vartheta+\vartheta\log(1+1/\vartheta)\le3$, $\vartheta H_2(\vartheta/\kappa)\le\vartheta^{1/2}\le\sqrt2$,
$\int_0^{K_f}H_1\le K_f+\int_0^{K_f}\log(1+\frac1y)\,dy=10+11\log11-10\log10<14$, and
$\int_0^{K_f}H_2\le2\sqrt{K_f}<7$. Hence
\begin{equation}\label{eq:Cstar}
E(t)-E(t_n)\le C_*vm_\lambda^2\lambda=C_*vM^2\lambda^{1+2\gamma},\qquad C_*:=d_LC_2\bigl(4+3+\sqrt2+21\kappa\bigr).
\end{equation}
$C_2$, and so $C_*$, involves only $c_6,C_\mu,C_T,\kappa,K_f,d_L$ and $1/m$: no $\eps$, $\lambda$, $R$ or $v$.
Finally $\frac12f_F^2+C_*v\le\frac18+\frac3{16}=\frac12\bar\theta$.
\end{proof}

\subsection{Length}\label{app:walk:length}

Stretching is driven by the same rates as bending, but it needs no cancellation: the length of a pair grows at rate at
most $v\int a^*(\rho)\,ds$, and the time the pursuer spends near the pair is short.

\begin{proposition}[Length]\label{prop:length}
If $v_0\le R/C_3$ for a suitable number $C_3$ (fixed in the proof), then $l_t(a,b)\le\frac98\lambda<\kappa\lambda$.
\end{proposition}

\begin{proof}
At the start time $l(t_\circ)\le(1+\delta)\lambda$, with $\delta:=0$ for body pairs and $\delta:=\delta_t\le\frac1{16}$ for
tail pairs (by (TW)). Always $l'=v\int T\cdot AT\,ds\le v\int a^*(\rho_s)\,ds$ (Lemmas~\ref{lem:Eprime}(i),
\ref{lem:F}(ii)). The pair is far on an initial segment of $[t_\circ,t]$ and near afterwards.

\emph{Far times.} There $\rho\ge\rhoarc\ge\eps$ on the arc (Proposition~\ref{prop:far}(i)), so $A=\Pi/\rho$ and
$(\log l)'\le v/\rhoarc$. (1) On $\{g>K_RR\}$ (duration $\le20$, $\rhoarc\ge\frac{47}8R$): increase of $\log l$ at most
$4v/R$. (2) On $\{c_{\mathrm{off}}R<g\le K_RR\}$, $(\log l)'\le4v/g$, and Lemma~\ref{lem:toolkit}(b)(i) with
$h(y)=y^{-1}\one\{y\le K_RR\}$, $g_{\min}>R/2$, gives $\int d\tau/g\le\frac R5\cdot\frac2R+\chi_g\log(2K_R)<9$: increase
at most $36v$. (3) On $\{g\le c_{\mathrm{off}}R\}$ use Lemma~\ref{lem:Eprime}(iv):
$(\log l)'\le(v/\rhoarc)\bigl[N/l^{1/2}+(|p_a|+l^{1/2}N/\pi)/\rhoarc\bigr]^2$. Here $\rhoarc\ge g/4$, $|p_a|\le C_pMg^{1+\gamma}$
(Proposition~\ref{prop:far}), $N/l^{1/2}\le M\lambda^{\gamma}\le Mg^\gamma$ and $l^{1/2}N\le\kappa^{1/2}M\lambda^{1+\gamma}\le Mg^{1+\gamma}\cdot\kappa^{1/2}$,
so $(\log l)'\le4(5+4C_p)^2vM^2g^{2\gamma-1}$. As $2\gamma-1>0$, Lemma~\ref{lem:toolkit}(b)(ii) with $\bar g=c_{\mathrm{off}}R$ and
$c_{\mathrm{off}}R+\eps/\kappa\le R$ gives $\int g^{2\gamma-1}d\tau\le\chi_gR^{2\gamma}/(2\gamma)\le\chi_gR^{2\gamma}$; the increase
is at most $4\chi_g(5+4C_p)^2vM^2R^{2\gamma}=4\chi_g(5+4C_p)^2m^2v$.

\emph{Near times.} Write $a^*(\rho)=\eps^{-1}\one\{\rho<\eps\}+\int_{\max(\rho,\eps)}^\infty y^{-2}dy$ (check: both
sides equal $1/\rho$ if $\rho\ge\eps$ and $2/\eps$ if $\rho<\eps$). With $\mathcal L(y):=$ arclength of
$\{s\in(a,b):\rho_s<y\}$, Fubini gives $l'\le v\bigl[\mathcal L(\eps)/\eps+\int_\eps^\infty\mathcal L(y)y^{-2}dy\bigr]$.
For $y\ge\eps$: if $\rho_s<y$, then $\Delta_s/c_6\le\max(\rho_s,\eps)\le y$ (Lemma~\ref{lem:toolkit}(c); $\Delta_s<K_RR$ as in the
envelope above), so $s-a\le c_6y$; the set lies in $(a,a+\min(\lambda,c_6y))$ and, by (J1) for this family pair
(length $\ge\eps/2$), $\mathcal L(y)\le\kappa\min(\lambda,c_6y)$.
(1) $\mathcal L(\eps)\ne0$ forces some $\Delta_s\le c_6\eps$, hence $g\le c_6\eps$; by Lemma~\ref{lem:toolkit}(b)(iii) this
happens on at most $\kappa c_6+1<7$ legs, each of duration $\le d_L\eps$, and there $\mathcal L(\eps)/\eps\le\kappa\lambda/\eps$.
Contribution at most $7d_L\kappa v\lambda$.
(2) Put $\bar\rho:=\max(\eps,\min_{[a,b]}\rho_s)$; $\mathcal L(y)=0$ for $\eps\le y<\bar\rho$, and
$c_6\bar\rho\ge\max(c_6\eps,g)$ since $\Delta_s\ge g$. With $h_1(z):=\int_z^\infty\min(\lambda,y)y^{-2}dy$ --- which equals
$1+\log(\lambda/z)$ for $z\le\lambda$ and $\lambda/z$ for $z>\lambda$, is non-increasing, satisfies $zh_1(z)\le\lambda$
and $\int_0^{K_f\lambda}h_1=\lambda(2+\log K_f)$ --- the substitution $z=c_6y$ gives
$\int_{\bar\rho}^\infty\mathcal L y^{-2}dy\le\kappa c_6h_1(c_6\bar\rho)\le\kappa c_6h_1(\max(c_6\eps,g))$. By
Lemma~\ref{lem:toolkit}(b)(i) with the non-increasing $h(y):=h_1(\max(c_6\eps,y))\one\{y<K_f\lambda\}$ (all near gaps are $<K_f\lambda$; and $h\le h_1$),
$\int\kappa c_6h_1(\max(c_6\eps,g(\tau)))\,d\tau\le\kappa c_6\bigl[d_L\eps\,h_1(c_6\eps)+\chi_g\int_0^{K_f\lambda}h_1\bigr]
\le\kappa\bigl[d_L+c_6\chi_g(2+\log10)\bigr]\lambda$.
So $l$ increases by at most $C_4v\lambda$ during the near phase, $C_4:=7d_L\kappa+\kappa d_L+\kappa c_6\chi_g(2+\log10)<81$.

\emph{Conclusion.} $l_t(a,b)\le\frac{17}{16}\lambda\exp\bigl(4v/R+36v+4\chi_g(5+4C_p)^2m^2v\bigr)+81v\lambda$, and the
right side is $\le\frac98\lambda$ once $v\le R/C_3$ for a suitable number $C_3$ (e.g.\ $C_3=2\cdot10^4$ works with
$m=\frac1{10}$, $C_p<9$).
\end{proof}

\subsection{Proof of the End Theorem}\label{app:walk:proof}

\begin{proof}[Proof of Theorem~\ref{thm:end}]
Take $v_0$ so small that \eqref{eq:V0} and the conditions of Propositions~\ref{prop:far}, \ref{prop:near},
\ref{prop:length} hold, and $\delta_t\le\frac1{16}$; all of them are of the form $v_0\le R/C$ or $v_0\le1/C$ with $C$
built from the locked parameters, so $v_0=v_0(m,R)$ and does not depend on $\eps$.

\emph{(a) and the proviso.} Let $t^*$ be as in Lemma~\ref{lem:cont}, $t\in[t_\td,t^*]$ and $(a,b)\in\FamS(t)$. If $(a,b)$
is a tail pair and $t=t_W$, (TW) gives $l\le\frac{17}{16}\lambda$ and $E\le\frac12s_0^2M^2\lambda^{1+2\gamma}$. Otherwise
Proposition~\ref{prop:length} gives $l\le\frac98\lambda$, and $E\le\frac12f_F^2M^2\lambda^{1+2\gamma}\le\frac12\bar\theta
M^2\lambda^{1+2\gamma}$ by Proposition~\ref{prop:far} if the pair is far at $t$, or $E\le\frac12\bar\theta M^2\lambda^{1+2\gamma}$
by Proposition~\ref{prop:near} if it is near (these propositions use (J1), (J2) only for pairs of $\FamS(\tau)$,
$\tau\le t$, by Lemma~\ref{lem:locality}). Thus (Imp) holds, and
Lemma~\ref{lem:cont} gives the proviso, the stop, (J1), (J2) on $[t_\td,T_\stopp]$, and that $t_W$ is reached if $\We>0$.
The chord bound in (a) is (CA) of Lemma~\ref{lem:toolkit}.

\emph{(b)} By (a), (J1) holds at every leg start, so Lemma~\ref{lem:toolkit}(a) applies with $t'=T_\stopp$: every
non-final leg moves the base by an amount in $[\eps/\kappa,\eps]$, there are at most $\kappa U/\eps$ of them, and each
lasts at most $\eps/(1-v)$ (Lemma~\ref{lem:chase}(1)); hence $T_\stopp\le t_\td+\kappa U/(1-v)$. The final leg $J+1$
has $a_{J+1}=U$, so the intervals $[a_j,a_{j+1}]$, $j\le J$, cover $[0,U]$, and by Lemma~\ref{lem:chase}(2) every true
label is seen by $T_\stopp$.

\emph{(c)} For $J\subset I$ in $\FamS(t)$, Lemma~\ref{lem:locality}(ii) (with $\tau=t$) shows that the pairs of the
family contained in $[a_{j(t)},\max(b,a_{j(t)+1})]$ are admissible, so they satisfy (J1), (J2) by (a), and
Lemma~\ref{lem:tiltchain} gives $|e_J-e_I|\le C_TM\lambda_I^\gamma$, with $C_T<6.11$ ($\gamma=\frac23$) resp.\
$C_T<5.58$ ($\gamma=\frac34$).

\emph{(d)} is Lemma~\ref{lem:toolkit}(c), with $c_6=(1+\delta_1)/c_A=\frac92$.

\emph{(e)} If $\tilde P=P$ on $[0,t_j]$, the fronts agree on $[t_\td,t_j]$ (Lemma~\ref{lem:chase}(4)), hence so do
the legs, bases, families and $t_W$ (if $t_W\le t_j$) up to $t_j$; (a), (c), (d) at times $t\le t_j$ involve only these
data.
\end{proof}
